% Options for packages loaded elsewhere
% Options for packages loaded elsewhere
\PassOptionsToPackage{unicode}{hyperref}
\PassOptionsToPackage{hyphens}{url}
\PassOptionsToPackage{dvipsnames,svgnames,x11names}{xcolor}
%
\documentclass[
  brazilian,
  letterpaper,
  DIV=11,
  numbers=noendperiod]{scrartcl}
\usepackage{xcolor}
\usepackage{amsmath,amssymb}
\setcounter{secnumdepth}{5}
\usepackage{iftex}
\ifPDFTeX
  \usepackage[T1]{fontenc}
  \usepackage[utf8]{inputenc}
  \usepackage{textcomp} % provide euro and other symbols
\else % if luatex or xetex
  \usepackage{unicode-math} % this also loads fontspec
  \defaultfontfeatures{Scale=MatchLowercase}
  \defaultfontfeatures[\rmfamily]{Ligatures=TeX,Scale=1}
\fi
\usepackage{lmodern}
\ifPDFTeX\else
  % xetex/luatex font selection
\fi
% Use upquote if available, for straight quotes in verbatim environments
\IfFileExists{upquote.sty}{\usepackage{upquote}}{}
\IfFileExists{microtype.sty}{% use microtype if available
  \usepackage[]{microtype}
  \UseMicrotypeSet[protrusion]{basicmath} % disable protrusion for tt fonts
}{}
\makeatletter
\@ifundefined{KOMAClassName}{% if non-KOMA class
  \IfFileExists{parskip.sty}{%
    \usepackage{parskip}
  }{% else
    \setlength{\parindent}{0pt}
    \setlength{\parskip}{6pt plus 2pt minus 1pt}}
}{% if KOMA class
  \KOMAoptions{parskip=half}}
\makeatother
% Make \paragraph and \subparagraph free-standing
\makeatletter
\ifx\paragraph\undefined\else
  \let\oldparagraph\paragraph
  \renewcommand{\paragraph}{
    \@ifstar
      \xxxParagraphStar
      \xxxParagraphNoStar
  }
  \newcommand{\xxxParagraphStar}[1]{\oldparagraph*{#1}\mbox{}}
  \newcommand{\xxxParagraphNoStar}[1]{\oldparagraph{#1}\mbox{}}
\fi
\ifx\subparagraph\undefined\else
  \let\oldsubparagraph\subparagraph
  \renewcommand{\subparagraph}{
    \@ifstar
      \xxxSubParagraphStar
      \xxxSubParagraphNoStar
  }
  \newcommand{\xxxSubParagraphStar}[1]{\oldsubparagraph*{#1}\mbox{}}
  \newcommand{\xxxSubParagraphNoStar}[1]{\oldsubparagraph{#1}\mbox{}}
\fi
\makeatother

\usepackage{color}
\usepackage{fancyvrb}
\newcommand{\VerbBar}{|}
\newcommand{\VERB}{\Verb[commandchars=\\\{\}]}
\DefineVerbatimEnvironment{Highlighting}{Verbatim}{commandchars=\\\{\}}
% Add ',fontsize=\small' for more characters per line
\usepackage{framed}
\definecolor{shadecolor}{RGB}{241,243,245}
\newenvironment{Shaded}{\begin{snugshade}}{\end{snugshade}}
\newcommand{\AlertTok}[1]{\textcolor[rgb]{0.68,0.00,0.00}{#1}}
\newcommand{\AnnotationTok}[1]{\textcolor[rgb]{0.37,0.37,0.37}{#1}}
\newcommand{\AttributeTok}[1]{\textcolor[rgb]{0.40,0.45,0.13}{#1}}
\newcommand{\BaseNTok}[1]{\textcolor[rgb]{0.68,0.00,0.00}{#1}}
\newcommand{\BuiltInTok}[1]{\textcolor[rgb]{0.00,0.23,0.31}{#1}}
\newcommand{\CharTok}[1]{\textcolor[rgb]{0.13,0.47,0.30}{#1}}
\newcommand{\CommentTok}[1]{\textcolor[rgb]{0.37,0.37,0.37}{#1}}
\newcommand{\CommentVarTok}[1]{\textcolor[rgb]{0.37,0.37,0.37}{\textit{#1}}}
\newcommand{\ConstantTok}[1]{\textcolor[rgb]{0.56,0.35,0.01}{#1}}
\newcommand{\ControlFlowTok}[1]{\textcolor[rgb]{0.00,0.23,0.31}{\textbf{#1}}}
\newcommand{\DataTypeTok}[1]{\textcolor[rgb]{0.68,0.00,0.00}{#1}}
\newcommand{\DecValTok}[1]{\textcolor[rgb]{0.68,0.00,0.00}{#1}}
\newcommand{\DocumentationTok}[1]{\textcolor[rgb]{0.37,0.37,0.37}{\textit{#1}}}
\newcommand{\ErrorTok}[1]{\textcolor[rgb]{0.68,0.00,0.00}{#1}}
\newcommand{\ExtensionTok}[1]{\textcolor[rgb]{0.00,0.23,0.31}{#1}}
\newcommand{\FloatTok}[1]{\textcolor[rgb]{0.68,0.00,0.00}{#1}}
\newcommand{\FunctionTok}[1]{\textcolor[rgb]{0.28,0.35,0.67}{#1}}
\newcommand{\ImportTok}[1]{\textcolor[rgb]{0.00,0.46,0.62}{#1}}
\newcommand{\InformationTok}[1]{\textcolor[rgb]{0.37,0.37,0.37}{#1}}
\newcommand{\KeywordTok}[1]{\textcolor[rgb]{0.00,0.23,0.31}{\textbf{#1}}}
\newcommand{\NormalTok}[1]{\textcolor[rgb]{0.00,0.23,0.31}{#1}}
\newcommand{\OperatorTok}[1]{\textcolor[rgb]{0.37,0.37,0.37}{#1}}
\newcommand{\OtherTok}[1]{\textcolor[rgb]{0.00,0.23,0.31}{#1}}
\newcommand{\PreprocessorTok}[1]{\textcolor[rgb]{0.68,0.00,0.00}{#1}}
\newcommand{\RegionMarkerTok}[1]{\textcolor[rgb]{0.00,0.23,0.31}{#1}}
\newcommand{\SpecialCharTok}[1]{\textcolor[rgb]{0.37,0.37,0.37}{#1}}
\newcommand{\SpecialStringTok}[1]{\textcolor[rgb]{0.13,0.47,0.30}{#1}}
\newcommand{\StringTok}[1]{\textcolor[rgb]{0.13,0.47,0.30}{#1}}
\newcommand{\VariableTok}[1]{\textcolor[rgb]{0.07,0.07,0.07}{#1}}
\newcommand{\VerbatimStringTok}[1]{\textcolor[rgb]{0.13,0.47,0.30}{#1}}
\newcommand{\WarningTok}[1]{\textcolor[rgb]{0.37,0.37,0.37}{\textit{#1}}}

\usepackage{longtable,booktabs,array}
\newcounter{none} % for unnumbered tables
\usepackage{calc} % for calculating minipage widths
% Correct order of tables after \paragraph or \subparagraph
\usepackage{etoolbox}
\makeatletter
\patchcmd\longtable{\par}{\if@noskipsec\mbox{}\fi\par}{}{}
\makeatother
% Allow footnotes in longtable head/foot
\IfFileExists{footnotehyper.sty}{\usepackage{footnotehyper}}{\usepackage{footnote}}
\makesavenoteenv{longtable}
\usepackage{graphicx}
\makeatletter
\newsavebox\pandoc@box
\newcommand*\pandocbounded[1]{% scales image to fit in text height/width
  \sbox\pandoc@box{#1}%
  \Gscale@div\@tempa{\textheight}{\dimexpr\ht\pandoc@box+\dp\pandoc@box\relax}%
  \Gscale@div\@tempb{\linewidth}{\wd\pandoc@box}%
  \ifdim\@tempb\p@<\@tempa\p@\let\@tempa\@tempb\fi% select the smaller of both
  \ifdim\@tempa\p@<\p@\scalebox{\@tempa}{\usebox\pandoc@box}%
  \else\usebox{\pandoc@box}%
  \fi%
}
% Set default figure placement to htbp
\def\fps@figure{htbp}
\makeatother



\ifLuaTeX
\usepackage[bidi=basic,shorthands=off]{babel}
\else
\usepackage[bidi=default,shorthands=off]{babel}
\fi
\ifLuaTeX
  \usepackage{selnolig} % disable illegal ligatures
\fi


\setlength{\emergencystretch}{3em} % prevent overfull lines

\providecommand{\tightlist}{%
  \setlength{\itemsep}{0pt}\setlength{\parskip}{0pt}}



 


\KOMAoption{captions}{tableheading}
\makeatletter
\@ifpackageloaded{caption}{}{\usepackage{caption}}
\AtBeginDocument{%
\ifdefined\contentsname
  \renewcommand*\contentsname{Índice}
\else
  \newcommand\contentsname{Índice}
\fi
\ifdefined\listfigurename
  \renewcommand*\listfigurename{Lista de Figuras}
\else
  \newcommand\listfigurename{Lista de Figuras}
\fi
\ifdefined\listtablename
  \renewcommand*\listtablename{Lista de Tabelas}
\else
  \newcommand\listtablename{Lista de Tabelas}
\fi
\ifdefined\figurename
  \renewcommand*\figurename{Figura}
\else
  \newcommand\figurename{Figura}
\fi
\ifdefined\tablename
  \renewcommand*\tablename{Tabela}
\else
  \newcommand\tablename{Tabela}
\fi
}
\@ifpackageloaded{float}{}{\usepackage{float}}
\floatstyle{ruled}
\@ifundefined{c@chapter}{\newfloat{codelisting}{h}{lop}}{\newfloat{codelisting}{h}{lop}[chapter]}
\floatname{codelisting}{Listagem}
\newcommand*\listoflistings{\listof{codelisting}{Lista de Listagens}}
\makeatother
\makeatletter
\makeatother
\makeatletter
\@ifpackageloaded{caption}{}{\usepackage{caption}}
\@ifpackageloaded{subcaption}{}{\usepackage{subcaption}}
\makeatother
\usepackage{bookmark}
\IfFileExists{xurl.sty}{\usepackage{xurl}}{} % add URL line breaks if available
\urlstyle{same}
\hypersetup{
  pdftitle={Aula Prática 3 -- Montagem Genômica e Avaliação de Qualidade},
  pdfauthor={Nicholas Barroso},
  pdflang={pt-BR},
  colorlinks=true,
  linkcolor={blue},
  filecolor={Maroon},
  citecolor={Blue},
  urlcolor={Blue},
  pdfcreator={LaTeX via pandoc}}


\title{Aula Prática 3 -- Montagem Genômica e Avaliação de Qualidade}
\author{Nicholas Barroso}
\date{2026-07-18}
\begin{document}
\maketitle

\renewcommand*\contentsname{Índice}
{
\setcounter{tocdepth}{3}
\tableofcontents
}

\section{Introdução}\label{introduuxe7uxe3o}

A montagem genômica permite reconstruir o genoma de um organismo a
partir de milhões de pequenos fragmentos de DNA (\emph{reads}) gerados
por tecnologias de sequenciamento de nova geração (NGS). O desafio está
em ordenar corretamente as leituras para obter o genoma completo.

A qualidade da montagem depende diretamente da qualidade dos dados de
entrada. Reads com baixa qualidade, presença de adaptadores ou
contaminações podem levar a montagens fragmentadas, com erros e regiões
não resolvidas. Por isso, o controle de qualidade e a limpeza (trimagem)
dos dados são etapas fundamentais antes de qualquer montagem.

Nesta aula prática faremos o fluxo de trabalho da montagem genômica,
desde a avaliação inicial dos dados até a comparação de diferentes
montadores (\emph{assemblers}). Utilizaremos as seguintes ferramentas:

\begin{itemize}
\tightlist
\item
  \textbf{FastQC}: para controle de qualidade dos reads;
\item
  \textbf{Trimmomatic}: para limpeza e remoção de adaptadores;
\item
  \textbf{SPAdes, Masurca e MEGAHIT}: três dos mais populares assemblers
  de genomas, cada um com abordagens algorítmicas distintas;
\item
  \textbf{QUAST}: para avaliação comparativa das montagens geradas.
\end{itemize}

\begin{center}\rule{0.5\linewidth}{0.5pt}\end{center}

\section{Objetivos}\label{objetivos}

\begin{itemize}
\tightlist
\item
  Realizar o controle de qualidade de dados de sequenciamento;
\item
  Executar a limpeza (trimagem) de reads;
\item
  Montar genomas utilizando diferentes montadores/\emph{assemblers}
  (SPAdes, Masurca e MEGAHIT);
\item
  Comparar e interpretar métricas de qualidade de montagens.
\end{itemize}

\begin{center}\rule{0.5\linewidth}{0.5pt}\end{center}

\section{Pré-requisitos}\label{pruxe9-requisitos}

\begin{itemize}
\tightlist
\item
  Conhecimento básico de linha de comando (Linux/Unix);
\item
  Ferramentas instaladas: FastQC, Trimmomatic, SPAdes, Masurca, MEGAHIT,
  QUAST;
\item
  Arquivos de entrada: \texttt{reads\_R1.fastq.gz} e
  \texttt{reads\_R2.fastq.gz} (dados paired-end);
\item
  Arquivo de adaptadores: \texttt{adapters.fa} (ex: TruSeq3-PE.fa).
\end{itemize}

\begin{center}\rule{0.5\linewidth}{0.5pt}\end{center}

\section{Fluxograma}\label{fluxograma}

\includegraphics[width=7.38in,height=31.18in]{pratica3_files/figure-latex/mermaid-figure-1.png}

\begin{center}\rule{0.5\linewidth}{0.5pt}\end{center}

\section{Resumo do Fluxo de Trabalho}\label{resumo-do-fluxo-de-trabalho}

\begin{verbatim}
Dados Brutos → QC → Limpeza → Montagem → Avaliação → Comparação → Conclusão
\end{verbatim}

\begin{center}\rule{0.5\linewidth}{0.5pt}\end{center}

\section{Parte 1 -- Controle de Qualidade com
FastQC}\label{parte-1-controle-de-qualidade-com-fastqc}

\subsection{Execução}\label{execuuxe7uxe3o}

Manual:
\href{https://www.bioinformatics.babraham.ac.uk/projects/fastqc/}{FastQC
Documentation}

\begin{Shaded}
\begin{Highlighting}[]
\ExtensionTok{fastqc}\NormalTok{ reads\_R1.fastq.gz reads\_R2.fastq.gz }\AttributeTok{{-}o}\NormalTok{ fastqc\_results/}
\end{Highlighting}
\end{Shaded}

\subsection{Interpretação dos
Relatórios}\label{interpretauxe7uxe3o-dos-relatuxf3rios}

Abra os arquivos HTML gerados e observe:

\begin{itemize}
\tightlist
\item
  \textbf{Qualidade por base (\emph{Per base sequence quality})}:
  Boxplots mostrando a pontuação Phred em cada posição. Regiões com
  qualidade \textless{} 20 indicam necessidade de trimagem.
\item
  \textbf{Conteúdo GC (\emph{Per sequence GC content})}: A distribuição
  deve ser próxima a uma curva normal. Picos anômalos podem indicar
  contaminação.
\item
  \textbf{Presença de adaptadores (\emph{Adapter Content})}: Identifica
  se adaptadores não foram totalmente removidos.
\item
  \textbf{Níveis de duplicação (\emph{Sequence Duplication Levels})}:
  Altos níveis podem indicar viés de PCR.
\end{itemize}

\emph{Com base nos relatórios, quais etapas de limpeza você considera
necessárias?}

\begin{center}\rule{0.5\linewidth}{0.5pt}\end{center}

\section{Parte 2 -- Limpeza (Trimagem) com
Trimmomatic}\label{parte-2-limpeza-trimagem-com-trimmomatic}

\subsection{Comando}\label{comando}

Trimmomatic é uma ferramenta flexível para remover adaptadores e bases
de baixa qualidade. O comando abaixo realiza a trimagem em dados
paired-end, gerando arquivos de reads limpos e descartando reads curtas.

\textbf{Manual}:
\href{http://www.usadellab.org/cms/?page=trimmomatic}{Trimmomatic
Manual}

\begin{Shaded}
\begin{Highlighting}[]
\ExtensionTok{TrimmomaticPE} \DataTypeTok{\textbackslash{}}
    \AttributeTok{{-}threads}\NormalTok{ 4 }\DataTypeTok{\textbackslash{}}
    \AttributeTok{{-}phred33} \DataTypeTok{\textbackslash{}}
\NormalTok{    reads\_R1.fastq.gz }\DataTypeTok{\textbackslash{}}
\NormalTok{    reads\_R2.fastq.gz }\DataTypeTok{\textbackslash{}}
\NormalTok{    forward\_paired.fastq.gz }\DataTypeTok{\textbackslash{}}
\NormalTok{    forward\_unpaired.fastq.gz }\DataTypeTok{\textbackslash{}}
\NormalTok{    reverse\_paired.fastq.gz }\DataTypeTok{\textbackslash{}}
\NormalTok{    reverse\_unpaired.fastq.gz }\DataTypeTok{\textbackslash{}}
\NormalTok{    ILLUMINACLIP:adapters.fa:2:30:10 }\DataTypeTok{\textbackslash{}}
\NormalTok{    LEADING:3 }\DataTypeTok{\textbackslash{}}
\NormalTok{    TRAILING:3 }\DataTypeTok{\textbackslash{}}
\NormalTok{    SLIDINGWINDOW:4:15 }\DataTypeTok{\textbackslash{}}
\NormalTok{    MINLEN:36}
\end{Highlighting}
\end{Shaded}

\subsection{Explicação dos
Parâmetros}\label{explicauxe7uxe3o-dos-paruxe2metros}

{\def\LTcaptype{none} % do not increment counter
\begin{longtable}[]{@{}
  >{\raggedright\arraybackslash}p{(\linewidth - 2\tabcolsep) * \real{0.5789}}
  >{\raggedright\arraybackslash}p{(\linewidth - 2\tabcolsep) * \real{0.4211}}@{}}
\toprule\noalign{}
\begin{minipage}[b]{\linewidth}\raggedright
Parâmetro
\end{minipage} & \begin{minipage}[b]{\linewidth}\raggedright
Função
\end{minipage} \\
\midrule\noalign{}
\endhead
\bottomrule\noalign{}
\endlastfoot
\texttt{ILLUMINACLIP:adapters.fa:2:30:10} & Remove adaptadores.
Parâmetros: mismatch máximo (2), palíndromo (30), simples (10) \\
\texttt{LEADING:3} & Remove bases com qualidade \textless{} 3 no início
da read \\
\texttt{TRAILING:3} & Remove bases com qualidade \textless{} 3 no final
da read \\
\texttt{SLIDINGWINDOW:4:15} & Corta a read quando a média de qualidade
em uma janela de 4 bases for \textless{} 15 \\
\texttt{MINLEN:36} & Descarta reads com comprimento final \textless{} 36
bases \\
\end{longtable}
}

\subsection{Verificação
Pós-Trimagem}\label{verificauxe7uxe3o-puxf3s-trimagem}

Execute novamente o FastQC nos arquivos \texttt{*\_paired.fastq.gz} para
confirmar a melhoria da qualidade.

\emph{Quantos reads foram descartados? Isso é esperado para seus dados?}

\begin{center}\rule{0.5\linewidth}{0.5pt}\end{center}

\section{Parte 3 -- Montagem com
SPAdes}\label{parte-3-montagem-com-spades}

Manual: \href{https://github.com/ablab/spades}{SPAdes GitHub}

\subsection{Execução}\label{execuuxe7uxe3o-1}

\begin{Shaded}
\begin{Highlighting}[]
\ExtensionTok{spades.py} \DataTypeTok{\textbackslash{}}
    \AttributeTok{{-}1}\NormalTok{ forward\_paired.fastq.gz }\DataTypeTok{\textbackslash{}}
    \AttributeTok{{-}2}\NormalTok{ reverse\_paired.fastq.gz }\DataTypeTok{\textbackslash{}}
    \AttributeTok{{-}o}\NormalTok{ montagem\_spades }\DataTypeTok{\textbackslash{}}
    \AttributeTok{{-}{-}careful} \DataTypeTok{\textbackslash{}}
    \AttributeTok{{-}t}\NormalTok{ 8 }\DataTypeTok{\textbackslash{}}
    \AttributeTok{{-}m}\NormalTok{ 64}
\end{Highlighting}
\end{Shaded}

\subsection{Parâmetros Importantes}\label{paruxe2metros-importantes}

\begin{itemize}
\tightlist
\item
  \texttt{-\/-careful}: Reduz erros em montagens de bactérias (modo mais
  rigoroso);
\item
  \texttt{-t\ 8}: Número de threads (ajuste conforme seu hardware);
\item
  \texttt{-m\ 64}: Limite de memória em GB.
\end{itemize}

\subsection{Saída Principal}\label{sauxedda-principal}

O arquivo de contigs está em: \texttt{montagem\_spades/contigs.fasta}

\begin{center}\rule{0.5\linewidth}{0.5pt}\end{center}

\section{Parte 4 -- Montagem com
Masurca}\label{parte-4-montagem-com-masurca}

Manual: \href{https://masurca.readthedocs.io/en/latest/}{Masurca
Documentation}

\subsection{Execução}\label{execuuxe7uxe3o-2}

\begin{Shaded}
\begin{Highlighting}[]
\ExtensionTok{masurca} \DataTypeTok{\textbackslash{}}
    \AttributeTok{{-}i} \StringTok{"forward\_paired.fastq.gz,reverse\_paired.fastq.gz"} \DataTypeTok{\textbackslash{}}
    \AttributeTok{{-}o}\NormalTok{ montagem\_masurca }\DataTypeTok{\textbackslash{}}
    \AttributeTok{{-}t}\NormalTok{ 8}
\end{Highlighting}
\end{Shaded}

\subsection{Processo}\label{processo}

O Masurca gera um script chamado \texttt{assemble.sh}. Execute-o para
iniciar a montagem:

\begin{Shaded}
\begin{Highlighting}[]
\BuiltInTok{cd}\NormalTok{ montagem\_masurca}
\ExtensionTok{./assemble.sh}
\end{Highlighting}
\end{Shaded}

\textbf{Importante:} Não use espaços após a vírgula na lista de
arquivos!

\subsection{Saída Principal}\label{sauxedda-principal-1}

O arquivo de contigs está em:
\texttt{montagem\_masurca/CA/primary.genome.scf.fasta}

\begin{center}\rule{0.5\linewidth}{0.5pt}\end{center}

\section{Parte 5 -- Montagem com
MEGAHIT}\label{parte-5-montagem-com-megahit}

Manual: \href{https://github.com/voutcn/megahit}{MEGAHIT GitHub}

\subsection{Execução}\label{execuuxe7uxe3o-3}

\begin{Shaded}
\begin{Highlighting}[]
\ExtensionTok{megahit} \DataTypeTok{\textbackslash{}}
    \AttributeTok{{-}1}\NormalTok{ forward\_paired.fastq.gz }\DataTypeTok{\textbackslash{}}
    \AttributeTok{{-}2}\NormalTok{ reverse\_paired.fastq.gz }\DataTypeTok{\textbackslash{}}
    \AttributeTok{{-}o}\NormalTok{ montagem\_megahit }\DataTypeTok{\textbackslash{}}
    \AttributeTok{{-}t}\NormalTok{ 8 }\DataTypeTok{\textbackslash{}}
    \AttributeTok{{-}m}\NormalTok{ 0.9}
\end{Highlighting}
\end{Shaded}

\subsection{Parâmetros}\label{paruxe2metros}

\begin{itemize}
\tightlist
\item
  \texttt{-m\ 0.9}: Usa 90\% da memória total disponível (ajuste
  conforme necessário);
\item
  \texttt{-t\ 8}: Número de threads.
\end{itemize}

\subsection{Saída Principal}\label{sauxedda-principal-2}

O arquivo de contigs está em:
\texttt{montagem\_megahit/final.contigs.fa}

\begin{center}\rule{0.5\linewidth}{0.5pt}\end{center}

\section{Parte 6 -- Avaliação Inicial das
Montagens}\label{parte-6-avaliauxe7uxe3o-inicial-das-montagens}

Para cada montagem, analise:

\begin{Shaded}
\begin{Highlighting}[]

\CommentTok{\# Contar número de contigs}
\FunctionTok{grep} \AttributeTok{{-}c} \StringTok{"\^{}\textgreater{}"}\NormalTok{ montagem\_spades/contigs.fasta}


\CommentTok{\# Tamanho total da montagem}
\FunctionTok{grep} \AttributeTok{{-}v} \StringTok{"\^{}\textgreater{}"}\NormalTok{ montagem\_spades/contigs.fasta }\KeywordTok{|} \FunctionTok{wc} \AttributeTok{{-}c}


\CommentTok{\# Maior contig}
\FunctionTok{grep} \AttributeTok{{-}v} \StringTok{"\^{}\textgreater{}"}\NormalTok{ montagem\_spades/contigs.fasta }\KeywordTok{|} \FunctionTok{awk} \StringTok{\textquotesingle{}\{print length($0)\}\textquotesingle{}} \KeywordTok{|} \FunctionTok{sort} \AttributeTok{{-}nr} \KeywordTok{|} \FunctionTok{head} \AttributeTok{{-}1}
\end{Highlighting}
\end{Shaded}

\textbf{Monte uma tabela comparativa:}

{\def\LTcaptype{none} % do not increment counter
\begin{longtable}[]{@{}llll@{}}
\toprule\noalign{}
Métrica & SPAdes & Masurca & MEGAHIT \\
\midrule\noalign{}
\endhead
\bottomrule\noalign{}
\endlastfoot
Nº de contigs & & & \\
Tamanho total (bp) & & & \\
Maior contig (bp) & & & \\
\end{longtable}
}

\begin{center}\rule{0.5\linewidth}{0.5pt}\end{center}

\section{Parte 7 -- Avaliação Detalhada com
QUAST}\label{parte-7-avaliauxe7uxe3o-detalhada-com-quast}

Manual: \href{http://quast.sourceforge.net/docs/manual.html}{QUAST
Documentation}

\subsection{Execução}\label{execuuxe7uxe3o-4}

\begin{Shaded}
\begin{Highlighting}[]
\ExtensionTok{quast.py} \DataTypeTok{\textbackslash{}}
    \AttributeTok{{-}o}\NormalTok{ quast\_results }\DataTypeTok{\textbackslash{}}
    \AttributeTok{{-}t}\NormalTok{ 8 }\DataTypeTok{\textbackslash{}}
    \AttributeTok{{-}m}\NormalTok{ 500 }\DataTypeTok{\textbackslash{}}
\NormalTok{    montagem\_spades/contigs.fasta }\DataTypeTok{\textbackslash{}}
\NormalTok{    montagem\_masurca/CA/primary.genome.scf.fasta }\DataTypeTok{\textbackslash{}}
\NormalTok{    montagem\_megahit/final.contigs.fa }\DataTypeTok{\textbackslash{}}
    \AttributeTok{{-}l} \StringTok{"SPAdes,Masurca,MEGAHIT"}
\end{Highlighting}
\end{Shaded}

\subsection{Métricas Chave a
Interpretar}\label{muxe9tricas-chave-a-interpretar}

{\def\LTcaptype{none} % do not increment counter
\begin{longtable}[]{@{}
  >{\raggedright\arraybackslash}p{(\linewidth - 4\tabcolsep) * \real{0.2368}}
  >{\raggedright\arraybackslash}p{(\linewidth - 4\tabcolsep) * \real{0.4211}}
  >{\raggedright\arraybackslash}p{(\linewidth - 4\tabcolsep) * \real{0.3421}}@{}}
\toprule\noalign{}
\begin{minipage}[b]{\linewidth}\raggedright
Métrica
\end{minipage} & \begin{minipage}[b]{\linewidth}\raggedright
O que significa
\end{minipage} & \begin{minipage}[b]{\linewidth}\raggedright
Valor ideal
\end{minipage} \\
\midrule\noalign{}
\endhead
\bottomrule\noalign{}
\endlastfoot
\textbf{N50} & Comprimento mínimo dos contigs que cobrem 50\% do genoma
& Quanto maior, melhor \\
\textbf{L50} & Número de contigs necessários para atingir o N50 & Quanto
menor, melhor \\
\textbf{Nº de contigs} & Fragmentação da montagem & Menor é melhor \\
\textbf{Maior contig} & Contig mais longo produzido & Quanto maior,
melhor \\
\textbf{GC (\%)} & Conteúdo GC estimado & Deve ser compatível com a
espécie \\
\textbf{Mismatches por 100 kbp} & Erros de base na montagem & Quanto
menor, melhor \\
\end{longtable}
}

\subsection{Análise Comparativa}\label{anuxe1lise-comparativa}

Abra o relatório HTML gerado (\texttt{quast\_results/report.html}) e
compare:

\begin{enumerate}
\def\labelenumi{\arabic{enumi}.}
\tightlist
\item
  Qual assembler produziu o melhor N50?
\item
  Qual teve o menor número de contigs?
\item
  Há diferenças significativas no conteúdo GC estimado?
\item
  Qual montagem você consideraria a melhor e por quê?
\end{enumerate}

\begin{center}\rule{0.5\linewidth}{0.5pt}\end{center}

\section{Parte 8 -- Perguntas para Discussão
Final}\label{parte-8-perguntas-para-discussuxe3o-final}

\begin{enumerate}
\def\labelenumi{\arabic{enumi}.}
\tightlist
\item
  \textbf{Qualidade dos dados}: Como a qualidade inicial dos reads
  afetou as montagens?
\item
  \textbf{Desempenho dos assemblers}: Quais foram as principais
  diferenças entre SPAdes, Masurca e MEGAHIT?
\item
  \textbf{Escolha do assembler}: Em quais situações você usaria cada um
  dos assemblers?
\item
  \textbf{Limitações}: Quais informações ainda faltam para validar
  completamente a montagem?
\item
  \textbf{Próximos passos}: O que você faria para melhorar ainda mais a
  montagem?
\end{enumerate}

\begin{center}\rule{0.5\linewidth}{0.5pt}\end{center}

\section{Referências e Leituras
Recomendadas}\label{referuxeancias-e-leituras-recomendadas}

\begin{itemize}
\tightlist
\item
  \href{https://www.bioinformatics.babraham.ac.uk/projects/fastqc/}{FastQC
  Documentation}
\item
  \href{http://www.usadellab.org/cms/?page=trimmomatic}{Trimmomatic
  Manual}
\item
  \href{https://github.com/ablab/spades}{SPAdes GitHub}
\item
  \href{https://masurca.readthedocs.io/en/latest/}{Masurca
  Documentation}
\item
  \href{https://github.com/voutcn/megahit}{MEGAHIT GitHub}
\item
  \href{http://quast.sourceforge.net/docs/manual.html}{QUAST
  Documentation}
\end{itemize}

\begin{center}\rule{0.5\linewidth}{0.5pt}\end{center}




\end{document}
